\subsection{I-29. A polynomial pair budget for arbitrarily many resets at one singular rank}
\label{I-29}

\paragraph{Statement and scope.}
The rank-profile mechanism used in I-21 has a growing-width
positive specialization.  Fix $1\le r<w$, put $c=w-r$, and
suppose a constant-boundary width-$w$ ROABP representation
of $f$ in one hidden variable order has the following promise:
every nonempty interval of its internal matrix word has generic
rank either $r$ or $w$.  The boundary row and column are not
included in this promise.  There may be arbitrarily many
rank-$r$ factors and intervals; their matrices need not commute.
Then the generator
\[
 X_i=t+\sum_{j=1}^{K}z_j\lambda_i(y_j),
 \qquad K=2r(w-r)+w-1,
 \tag{I.42}
\]
hits every nonzero such $f$ in every characteristic, over a
field or extension supplying the requisite distinct nodes.
For total degree at most $D$, the explicit grid construction
of I-21 has size at most
\[
 (D+1)\bigl((n-1)D+1\bigr)^K\binom{D+K}{K}.
 \tag{I.43}
\]
Its field-size convention is the same as I-21.

This is a polynomial \emph{pair budget} in $w$, with
sampling cost $(nD)^{O(w^2)}$, not polynomial-size PIT for
variable width.  The statement permits unbounded total
corank and so is different from I-20.  It further specializes
the cited unit-rank mechanism; no historical novelty is claimed.

\paragraph{1. A two-rank coefficient-space lemma.}
Let $B_i(Y_i)$ be $w\times w$ polynomial matrices in distinct
variables over a field $L$.  Assume every interval has generic
rank equal to its rank at zero, and that these ranks belong
to $\{r,w\}$.  Then, for every constant row $b$,
\[
 \begin{gathered}
 \operatorname{Coeff}(bB_1\cdots B_m)
 =\operatorname{span}_L\{
       [Y^\alpha](bB_1\cdots B_m):
       |\operatorname{supp}\alpha|\le h\},\\
 h=r(w-r)+w-1=rc+w-1.
 \end{gathered}
 \tag{I.44}
\]
The assertion remains valid in a truncated product local ring
when the corresponding interval unit-rank assumptions hold.

If the full constant product is invertible, all constant
factors are invertible, and I-17's nested row-space argument
gives the smaller bound $w-1$.  Otherwise its rank is $r$.
Choose a persistent $r$-dimensional row space $U_0$ on which
the full constant product is injective and transport it to
$U_i$ at every cut.  Choose a complement $V_m$ of $U_m$;
the suffix pullback
\[
 V_i=\{z:zB_{i+1}(0)\cdots B_m(0)\in V_m\}
\]
has dimension $c$, complements $U_i$, and is carried into
$V_{i+1}$ by the next constant matrix.  In compatible bases
the constants therefore have the form
\[
 B_i(0)=\operatorname{diag}(I_r,Q_i),
 \tag{I.45}
\]
where $Q_i$ is either invertible or zero: the rank of the
whole factor is either $w$ or $r$.  These basis changes
preserve supports and interval ranks.

Let $\pi:L^{1\times w}\to L^{1\times r}$ be active-coordinate
projection.  Prefix coefficient spaces $S_i$ have nested
projections $\pi S_i$, by (I.45).
We need the following stable-segment assertion: on a segment
where every projected space is the same $T$, give basis lifts
of $T$ initial cost zero and a basis of the initial projection
kernel cost one.  Cost means initial cost plus the number of
nonconstant local coefficients used.  Then cost at most $c$
spans the final coefficient space.

First suppose all $Q_i$ in the segment are invertible.
After constant transport, the coefficient spaces are nested
inside $\pi^{-1}(T)$, whose dimension is $\dim T+c$.
If the initial kernel has dimension $k>0$, all initial rows
cost at most one and there can be at most $c-k$ subsequent
dimension increases.  At each increase one nonconstant
coefficient suffices, so the cost is at most
$1+c-k\le c$.  If $k=0$, the initial cost is zero and
there are at most $c$ increases.  Stable layers require only
constant propagation.  This proves the assertion in this case.

If a zero $Q_i$ occurs, take its last position $a$ and put
$P=B_a\cdots B_b$, where $b$ is the segment's final position.
The suffix has rank $r$ and constant matrix
$\operatorname{diag}(I_r,0)$: its rank is at most $r$ from
the factor at $a$, and its constant $r$-minor is nonzero.
In the power-series ring at zero, or in the truncated local
ring, write its top blocks as $[U\ V]$.  The matrix $U$
is invertible.  The vanishing $(r+1)$-minors give the row-graph
identity
\[
 (x,y)\in\operatorname{rowmod}(P)
       \quad\Longrightarrow\quad y=xU^{-1}V.
\]
Let $B$ be the matrix of the initial lifts propagated through
the constants before position $a$, and let $H=\pi B$ be the
chosen basis of $T$.  Stable projection implies
\[
 \pi(BP)=AH,\qquad \pi(sP)=\eta H\quad(s\in S_{a-1}),
 \qquad A(0)=I.
\]
Thus $A$ is invertible and the row graph gives
$sP=\eta A^{-1}BP$.  If $T=0$, it instead gives $sP=0$.
It follows, by taking coefficient rows, that the final space
is exactly $\operatorname{Coeff}(BP)$.  Indeed, the displayed
identity proves one containment, and the rows of $B$ already
belong to $S_{a-1}$ for the other.

Restart after position $a$ from
$\operatorname{Coeff}(BB_a)$.  It has constant lifts of $T$
of cost zero and a kernel basis of cost at most one, since
only that one variable position was used.  The remaining
constants are invertible.  The preceding dimension argument
then bounds the total cost by $c$, proving the stable-segment
assertion.  In the untruncated case the inverses may be formal
power series; each coefficient of their product with $BP$
is still a finite linear combination of coefficient rows of
$BP$.  No analytic convergence is used.

We now count support from the original single row $b$.
If $\pi b\ne0$, there are at most $r-1$ strict projection
growths.  The initial stable segment costs at most $c$;
each growth factor and the following stable segment add at
most $1+c$.  The bound is
\[
 c+(r-1)(c+1)=rc+r-1\le h.
\]
If $\pi b=0$, the first stable segment lies in the passive
$c$-space.  If it contains a zero $Q_i$, the row-graph
argument kills the entire current coefficient space and the
final word is zero.  Otherwise its constants are invertible,
and starting from one row there are at most $c-1$ dimension
increases.  This first segment costs at most $c-1$.
At most $r$ projection growths follow, giving the bound
\[
 (c-1)+r(c+1)=rc+r+c-1=rc+w-1=h.
\]
At every stage the retained rows are linear combinations of
actual coefficient rows of the original prefix with the
stated support.  At a restart the chosen lifts are propagated
by constants through skipped positions; selecting a local
nonconstant coefficient adds only its own variable position.
Thus the cost count proves the actual support assertion
(I.44), rather than just a bound on an abstract generating
module.  The same reasoning applies in the truncated ring.

\paragraph{2. Sharpness for every active and passive dimension.}
Write the active coordinates as $a_1,\ldots,a_r$ and the
passive coordinates as $p_1,\ldots,p_c$, and put
$P=\operatorname{diag}(I_r,0_c)$.
Use a fresh variable in every factor below.
First apply the $c-1$ unipotent chain factors
\[
 I+Y_1E_{p_1,p_2},\ I+Y_2E_{p_2,p_3},\ \ldots,
 I+Y_{c-1}E_{p_{c-1},p_c}.
\]
This initial list is empty if $c=1$.
For each $j=1,\ldots,r$, apply one reset
\[
 P+S_jE_{p_c,a_j},
\]
followed by $c$ unipotent chain factors for the edges
$a_j\to p_1\to p_2\to\cdots\to p_c$.
The total number of factors is
$(c-1)+r(c+1)=h$.

Each reset has generic rank exactly $r$: its only extra
row is a multiple of the existing active row at $a_j$.
Each chain factor is unipotent and invertible.  Every
interval without a reset has constant matrix $I_w$ and
generic rank $w$.  Every interval with a reset has constant
matrix $P$, has generic rank at most $r$ from the reset
factor, and has rank at least $r$ from its constant active
minor.  Thus every interval satisfies precisely the
two-rank promise, including rank preservation at zero.

Starting with $e_{p_1}^T$, the coefficient in the last
passive coordinate after the initial chain is the product
of all its $c-1$ variables.  At reset $j$, all passive
coordinates are killed except that last coefficient,
which is transferred, multiplied by $S_j$, into the
fresh active coordinate $a_j$.  Earlier active coordinates
are unchanged.  The following chain multiplies that
coefficient by each of its $c$ fresh variables before
placing it in $p_c$.  The earlier active coordinates
do not enter this chain.  Induction on $j$ proves that
the final contraction with $e_{p_c}$ is exactly the
product of all $h$ variables.

Consequently every coefficient row of support less than
$h$ is annihilated by $e_{p_c}$, whereas the full-support
coefficient is not.  This proves sharpness of (I.44) for
every $1\le r<w$.  In particular the corank-one case has
$h=2(w-1)$, and balanced active/passive dimensions require
quadratic support.  The scalar witness is a monomial and
is hit by one generic dense direction; this is not a
lower bound on dense-generator pair count.

\paragraph{3. Uniform preservation and the sampler.}
Start at the common transcendental center $X_i=t$.
All rank-$w$ intervals retain rank: each of their univariate
factors has a nonzero determinant, and the product of those
determinants remains nonzero at $t$.

For every rank-$r$ interval, use I-21's finite local-rank-test
construction.  Minimality makes each test's interior chunk
have rank greater than $r$, and the two-rank promise makes
that rank $w$.  Its selected nonzero $r$-minor is therefore
a contraction of an invertible compound chunk with at most
two endpoint factors.  Its endpoint compound row is
decomposable, and both endpoint vectors are nonzero at the
common transcendental center, exactly as checked in I-28.
That result gives a support budget $g=r(w-r)=rc$ for the
whole test, including the endpoints.

Add $g$ fresh SV pairs.  Each selected nonzero minor stays
nonzero, and their finite product is nonzero, so all such
tests and hence all rank-$r$ intervals are simultaneously
preserved over the parameter field.  Rank-$w$ intervals
remain preserved because the new amplitudes may be set to
zero.  This is an existence argument over one common
parameter field and requires neither the hidden order nor
the selected minors as input to the generator.

At the resulting center, translate all original variables
by independent residual variables.  Every interval now
has its generic rank at zero, in $\{r,w\}$.  The original
scalar remains a nonzero polynomial, since translation
over that field is invertible.  Apply (I.44) to its starting
row and then its final column.  Some scalar monomial has
support at most $h$; otherwise the final column would
annihilate the spanning coefficient rows and the whole
polynomial would be zero.  Adding $h$ further fresh SV pairs
therefore makes the scalar nonzero.  The two additive maps
combine into $K=g+h=2rc+w-1$ pairs, proving (I.42).

For sampling, the degree in $t$ is at most $D$, each selector
degree is at most $(n-1)D$, and the total amplitude degree
is at most $D$.  The same product grid and amplitude Newton
simplex as in I-21 give (I.43).  All arguments and the
interpolation convention are characteristic independent.
\hfill$\square$

\paragraph{What remains outside this positive case.}
I-24's exponentially long support witnesses use several
different singular ranks, rather than only one fixed rank
and full rank.  In that setting the passive blocks in
(I.45) may themselves have intermediate nonzero singular
ranks.  The last-zero restart used above then does not
replace the whole passive block, and recursive rank losses
can accumulate.  I-29 identifies a precisely proved regime
where this recursion is unnecessary.  It neither asserts
that recursion is necessary for dense nonvanishing nor
excludes rational-function or cocharacter constructions
for the remaining profiles.
